\chapter{Вещественный KO6-подъём коэффициентного цикла Тёплица}

Комплексная граничная карта дала два сопряжённых оператора с индексами
$-15$ и $+15$. Настоящий гейт проверяет, можно ли собрать из них один
вещественный KO6-цикл и сохранить ненулевой индекс полного пакета.

Проверка опирается на первичную литературу по вещественным
Kasparov-модулям и real bulk--edge correspondence. Бурн, Келлендонк и
Ренни подчёркивают, что антилинейные симметрии требуют $KKO$, а не
простого комплексного $KK$; см. \path{arXiv:1604.02337}. Рубин и
Домбровский показывают, что продолжение вещественной структуры на
расширенный спектральный объект не единственно и его KO-размерность должна
вычисляться для конкретного $J$; см. \path{arXiv:2012.15698}. Общий
современный справочник по вещественной $K$-теории --- версия 4 работы
Бурсема и Шокета от 29 июля 2026 года,
\path{arXiv:2407.05880v4}.

\section{Знаки текущего KO6}

В замороженной конечной геометрии используются
\begin{equation}
 J^2=1,
 \qquad
 JD=DJ,
 \qquad
 J\gamma=-\gamma J.
 \label{eq:v5-real-toeplitz-ko6-signs}
\end{equation}
Последнее равенство означает, что $J$ меняет две градуированные половины
местами.

Пусть нечётный самосопряжённый фредгольмов модуль имеет вид
\begin{equation}
 F=
 \begin{pmatrix}
  0&T^*\\
  T&0
 \end{pmatrix},
 \qquad
 \gamma=
 \begin{pmatrix}
  I&0\\0&-I
 \end{pmatrix}.
 \label{eq:v5-real-toeplitz-F}
\end{equation}

\section{Теорема о компенсации индекса}

Если $x\in\ker T$, то $x$ лежит в положительной градуированной половине
и $Fx=0$. Из $JF=FJ$ следует
\begin{equation}
 F(Jx)=J(Fx)=0.
\end{equation}
Поскольку $J\gamma=-\gamma J$, вектор $Jx$ лежит в отрицательной
половине, а значит принадлежит $\ker T^*$. Обратное отображение задаётся
$J^{-1}=J$.

\begin{theorem}[KO6-компенсация обычного индекса]
Для всякого фредгольмова цикла, удовлетворяющего
$J^2=1$, $JF=FJ$ и $J\gamma=-\gamma J$, вещественная структура задаёт
антилинейную биекцию
\begin{equation}
 J:\ker T\xrightarrow{\simeq}\ker T^*.
\end{equation}
Поэтому
\begin{equation}
 \operatorname{ind}T
 =\dim\ker T-\dim\ker T^*=0.
 \label{eq:v5-real-toeplitz-zero-index-theorem}
\end{equation}
\end{theorem}

Следовательно, ненулевой комплексный индекс одной ориентированной ветви
не может одновременно быть обычным индексом полного KO6-пакета. Это
структурное следствие знаков, а не недостаток выбранного оператора
Тёплица.

\section{Сбалансированный цикл}

Обозначим
\begin{equation}
 T_+=S\otimes q_0,
 \qquad
 T_-=S^*\otimes\overline{q_0}.
\end{equation}
На удвоенном градуированном модуле положим
\begin{equation}
 T_{\mathbb R}=T_+\oplus T_-.
 \label{eq:v5-real-toeplitz-balanced-T}
\end{equation}
Тогда
\begin{align}
 \dim\ker T_{\mathbb R}&=15,\\
 \dim\ker T_{\mathbb R}^*&=15,\\
 \operatorname{ind}T_{\mathbb R}&=0.
\end{align}
Антилинейный оператор $J$ комплексно сопрягает коэффициенты, меняет две
ориентации и две градуированные половины местами. Его фазы можно выбрать
так, что выполняются все три знака
\eqref{eq:v5-real-toeplitz-ko6-signs}.

Компактный дефект при этом не исчезает. В одной половине он расположен в
коядре $S$, в другой --- в ядре $S^*$. Суммарный компактный ранг равен
тридцати, а коэффициентный носитель удваивается со $105$ до $210$:
\begin{equation}
 \frac{30}{210}=\frac17.
 \label{eq:v5-real-toeplitz-weight-survives}
\end{equation}

\begin{proposition}[Вес сохраняется, индекс компенсируется]
KO6-удвоение сохраняет нормированный дефектный вес $1/7$, но заставляет
обычные целочисленные индексы двух ориентаций взаимно сократиться.
Ориентированный индекс $\pm15$ относится к комплексной половине после
забывания $J$, а не к полному вещественному пакету.
\end{proposition}

\section{Статус вещественного класса}

Нулевой обычный индекс не означает автоматически тривиальность во всех
вещественных теориях. Real $K$-теория различает выбранную вещественную
форму алгебры, инволюцию пространства и Clifford-степень. Бурн строит
инварианты симметричных квазисвободных состояний в группах
$KO_*(A^{\mathfrak r})$ через Clifford-модульные индексы; см.
\path{arXiv:2108.01230}. В физических применениях real Toeplitz extensions
могут иметь также $\mathbb Z_2$-граничные классы, но конкретная группа
зависит от класса симметрии и Real-инволюции.

Для текущего проекта ещё не зафиксировано, какой именно вещественной
$C^*$-алгеброй является коэффициентный объект:
\begin{enumerate}
 \item $M_{105}(\mathbb C)$, забытый до вещественной алгебры;
 \item фиксированная алгебра относительно $J$;
 \item Real-алгебра с инволюцией, обращающей хопфову окружность;
 \item Clifford-стабилизированная версия одной из них.
\end{enumerate}
Без этого нельзя честно назвать группу $KO_j$ или $KR_j$, в которой мог
бы жить остаточный класс. Последняя версия справочника
\path{arXiv:2407.05880v4} особо подчёркивает роль коммутативных вещественных
алгебр, не сводящихся к $C(X,\mathbb R)$, и различие между $KO$, $KU$ и
$KR$.

\begin{nogocriterion}[Запрет угадывания вещественной группы]
Из одной комплексной пары индексов $-15/+15$ нельзя заключить наличие
ненулевого $\mathbb Z_2$- или целочисленного Real-класса. Сначала должны
быть объявлены вещественная алгебра, её инволюция, Clifford-степень и
граничная карта соответствующей точной последовательности.
\end{nogocriterion}

\section{Что получено и что осталось}

Сбалансированный ограниченный вещественный фредгольмов цикл существует и
совместим со знаками текущего KO6. Он не вводит асимметрии между частицей
и сопряжённой частицей и сохраняет следовой вес $1/7$.

Однако полный ненулевой целочисленный индекс запрещён самой вещественной
структурой. Поэтому следующий вопрос меняется: нужно не пытаться ещё раз
получить индекс пятнадцать у полного пакета, а классифицировать его как
Real/KR-класс и определить, остаётся ли после компенсации нетривиальный
симметрийно-защищённый остаток.

\begin{protocol}
\begin{enumerate}
 \item знаки $J^2=1$, $JF=FJ$, $J\gamma=-\gamma J$:
 \textbf{реализованы};
 \item вещественный сбалансированный цикл Тёплица: \textbf{построен};
 \item индексы ориентированных половин: \textbf{$-15/+15$};
 \item обычный индекс полного KO6-пакета: \textbf{ноль};
 \item нулевой индекс является случайной компенсацией: \textbf{нет,
 следует из KO6};
 \item суммарный дефектный вес: \textbf{$1/7$};
 \item конкретная группа $KO/KR$: \textbf{не зафиксирована};
 \item ненулевой Real-класс: \textbf{не доказан};
 \item неограниченный родительский оператор: \textbf{не построен};
 \item физическая локализация и энергия: \textbf{не получены};
 \item следующий гейт: \textbf{классификация Real/KR-класса}.
\end{enumerate}
\end{protocol}

Машинный сертификат сохранён в
\path{s2t_v5_real_toeplitz_ko6_parent_lift_gate_results.json}.